% ************************** Thesis Abstract *****************************
% Use `abstract' as an option in the document class to print only the titlepage and the abstract.
\begin{abstract}
The modernization of the University of San Carlos (USC) Clinic’s Patient Information System (PIS) aims to resolve significant operational inefficiencies caused by legacy paper-based workflows. The previous system struggled with fragmented health records, slow retrieval times, and limited capabilities for health information dissemination and data-driven reporting. This project deployed a robust, three-tier web-based architecture utilizing React 18, Django 5.0.2, and PostgreSQL, specifically designed to modernize clinical record management and medical certificate issuance (USC Form ACA-HSD-04F). To ensure strict compliance with data privacy regulations, the system integrates the \texttt{pgcrypto} extension, enforcing column-level symmetric encryption (AES-256) for all Protected Health Information (PHI).

Rigorous Software Quality Assurance (SQA) audits and performance benchmarking demonstrated exceptional system reliability. The optimized database architecture successfully handled search queries with an average latency of 127.81ms, operating nearly four times faster than the 500ms maximum acceptable threshold. Furthermore, the custom PDF rendering module generated official medical certificates, optimally rendering between 122.93ms and 124.53ms. During the primary User Acceptance Testing (UAT), the student demographic (N=70) reported an overall satisfaction mean of 4.76 out of 5.00, while clinic staff (N=7) validated the system's operational readiness. 

To explicitly measure the digital transition's impact, a secondary comparative simulation was conducted involving students (N=30) and clinic staff (N=8) with prior experience in the legacy manual processes. This simulation yielded overall satisfaction means of 4.73 and 4.17 out of 5.00, respectively, and mathematically proved the system's superiority over the paper-based baseline. The empirical data demonstrated massive, quantifiable operational upgrades in medical record retrieval, health campaign dissemination, and medical certificate issuance speed. Ultimately, this rigorous comparative testing definitively justified the successful attainment of all core project objectives, establishing the USC-PIS as a comprehensive and matured solution for university healthcare delivery.
\end{abstract}
